\subsection{Consequences of the proved influence bound}
\label{ian:influence-form-consequences}

Status: proved for the listed exchangeable, Curie--Weiss, independent-block, and total-interaction-budget conclusions, using Corollary~\ref{ian:product-influence-aa-proved} of Subsection~\ref{ian:influence-form-uniform-aa} (no (UAA)) and, for the exchangeable laws, the slice-transfer theorem of Girish--Mittal--Nehoran--Raz cited from the literature; this subsection records the recomputed constants now used in Subsections~\ref{ian:exchangeable-generalized-influence}, \ref{ian:generalized-influence-checks}, \ref{ian:gamma-set-blocks-and-forests} and~\ref{ian:gamma-set-high-temperature}. The unrestricted weighted-row-sum Gamma question remains unresolved.

Throughout this subsection,
\[
 c_\star=\frac{9}{2^{61}\pi^{12}}.
\]
All polynomial bounds retain the entire-cube boundedness hypothesis.

\begin{proposition}[Updated product, slice, and exchangeable constants]
\label{ian:influence-consequences-exchangeable}
In \path{ian/04_exchangeable_generalized_influence.tex}, the product
bound is supplied by Corollary~\ref{ian:product-influence-aa-proved}
and the constants are
\[
 \kappa=c_\star,\qquad
 \lambda=10^{-9}2^{-146}c_\star^4,\qquad
 \sigma=2^{-112}\lambda.
 \tag{EC}
\]
With these, Lemma~\ref{ian:uniform-product-slice-aa} holds without
assuming (UAA), with exactly its stated exponents:
\[
 \operatorname{Var}_{\rho_{n,p}}f\ge t
 \ \Longrightarrow\
 \max_i I_i^{\rho_{n,p}}(f)\ge\kappa(t/d)^{12},
\]
\[
 \operatorname{Var}_{\omega_{n,k}}f\ge t
 \ \Longrightarrow\
 \max_i I_i^{\omega_{n,k}}(f)\ge\lambda(t/d)^{56}.
\]
Theorem~\ref{ian:exchangeable-gamma} consequently holds with
\[
 c_L=\frac{\sigma^2}{264}
     =\frac{10^{-18}2^{-516}c_\star^8}{264},
 \qquad C_L=114,
\]
\[
 c_G=\frac1{18},\qquad C_G=1,\qquad B_G=128,\qquad K_G=4.
\]
These constants are uniform over all exchangeable laws and their
positive-probability pinnings.
\end{proposition}

\begin{proof}
The product assertion follows directly from (P-AA), so the earlier
constants \(c_0,A_0\) and the product encoding are unnecessary.

The slice step retains the literature input already used in
Lemma~\ref{ian:uniform-product-slice-aa}:
Girish--Mittal--Nehoran--Raz, \emph{Distributional Variants of the
Aaronson--Ambainis Conjecture}, arXiv:2609.35327v1,
Proposition~4.12(1), PDF pages~20--21.
Its quantitative transfer is
\[
 S'(u,d)\ge10^{-9}d^{-8}
             \min\{u^5,\ uS(u/2,d)^4\}.
\]
This is a citable slice-transfer theorem, not an additional AA
hypothesis. Its proof has no density-dependent constant.
The slice continuity input remains its Lemma~4.9, PDF page~19,
as recorded in Lemma~\ref{ian:slice-continuity}.
% [Girish--Mittal--Nehoran--Raz, Proposition 4.12 and Lemma 4.9](https://arxiv.org/pdf/2609.35327v1#page=19)

For \(g(y)=(1+f(2y-1))/2\), variance is divided by four and
unnormalized flip influence equals \(I_i(f)\).
The product theorem supplies the conservative bound
\(S(u,d)=\kappa(u/d)^{12}\) in those conventions. Applying the
transfer with \(u=t/4\) gives
\[
 \max_i I_i^{\omega_{n,k}}(f)
 \ge10^{-9}d^{-8}
 \min\left\{(t/4)^5,\
       (t/4)\kappa^4\left(\frac{t}{8d}\right)^{48}\right\}
 =\lambda t^{49}d^{-56}
 \ge\lambda(t/d)^{56}.
\]
The minimum is the second term: \(\kappa=c_\star\le1\) and
\(t/(8d)\le t/4\le1\), so
\(\kappa^4(t/(8d))^{48}\le(t/4)^{48}\le(t/4)^4\).
Its value is \(10^{-9}2^{-2}2^{-144}\kappa^4t^{49}d^{-56}\), and
\(t^{49}\ge t^{56}\) since \(t\le1\).
The endpoint slices have zero variance.

The remaining proof of Theorem~\ref{ian:exchangeable-gamma} is
unchanged. Its within-slice threshold is
\[
 s_*=\lambda(v/(4d))^{56}=\sigma(v/d)^{56},
\]
and its localization step gives
\[
 \max_i I_i^\mu(f)
 \ge\frac{v s_*^2}{264d^2}
 \ge\frac{\sigma^2}{264}(v/d)^{114}.
\]
The between-slice argument retains all four Global parameters.
\end{proof}

\begin{remark}[The degree-\(db\) encoding route is not used]
\label{ian:influence-consequences-dyadic}
The product part of the proof in
\path{ian/04_exchangeable_generalized_influence.tex} now uses
(P-AA) directly. It is not a verbatim substitution of the new uniform
theorem into the earlier dyadic-encoding calculation, which would
lose exponents.

If only Theorem~\ref{ian:uniform-influence-aa-proved} were used,
the encoded polynomial \(h=f\circ g\) would have degree at most
\(db\), and the new source estimate would be
\[
 \max_{i,j}I_{ij}^{U_{nb}}(h)
 \ge\frac{c_\star(t/2)^{12}}{d^{12}b^{12}}.
\]
With \(x=d/t\ge1\), choose an integer \(A\ge64\) satisfying
\[
 2^{16}(A+1)^{12}2^{-A}\le c_\star,
 \qquad b=\lceil Ax\rceil,
 \qquad
 L=\frac{c_\star t^{12}}{2^{13}d^{12}b^{12}}.
\]
For the original continuity error \(\varepsilon=8d^22^{-b}\),
\[
 \varepsilon/t\le8x^22^{-Ax}\le1/2,\qquad
 \varepsilon/L
 \le2^{16}(A+1)^{12}c_\star^{-1}x^{26}2^{-Ax}\le1.
\]
Both right sides are nonincreasing on \(x\ge1\) because
\(A>26/\ln2\approx37.5\), which \(A\ge64\) guarantees; at \(x=1\)
they are at most \(8\cdot2^{-64}\le1/2\) and, by the choice of \(A\),
at most \(1\). The earlier floor \(A_0\ge32\) would not suffice for
this monotonicity, since \(32\ln2\approx22.2<26\).
Thus that route gives only
\[
 \max_i I_i^{\rho_{n,p}}(f)
 \ge\kappa_{\mathrm{enc}}(t/d)^{24},
 \qquad
 \kappa_{\mathrm{enc}}
 =\frac{c_\star}{2^{13}(A+1)^{12}}.
\]
Sketch, not separately checked: propagating this exponent through
the same slice and localization computations, with \(12\) replaced
by \(24\), gives
\[
 \lambda_{\mathrm{enc}}
 =10^{-9}2^{-290}\kappa_{\mathrm{enc}}^4,
 \qquad
 \max_iI_i^{\omega_{n,k}}(f)
 \ge\lambda_{\mathrm{enc}}(t/d)^{104},
\]
\[
 \sigma_{\mathrm{enc}}=2^{-208}\lambda_{\mathrm{enc}},
 \qquad
 c_{L,\mathrm{enc}}=\sigma_{\mathrm{enc}}^2/264,
 \qquad C_{L,\mathrm{enc}}=210.
\]
The direct product theorem avoids these losses and gives the
constants in (EC), with \(C_L=114\); nothing below depends on this
remark.
\end{remark}

\begin{corollary}[The Curie--Weiss check]
\label{ian:influence-consequences-cw}
Proposition~\ref{ian:cw-generalized-check} in
\path{ian/05_generalized_influence_checks.tex} holds without
assuming (UAA), with the six constants of
Proposition~\ref{ian:influence-consequences-exchangeable}.
Its proof is unchanged after replacing the invocation of
Theorem~\ref{ian:exchangeable-gamma} by its updated version.

The low-temperature failure of a standalone Local theorem, the
magnetisation reveal calculation, and the character-code refutation
retain their original statements and proofs.
\end{corollary}

\begin{proposition}[Updated block resampling and two-spin bounds]
\label{ian:influence-consequences-blocks}
Lemma~\ref{ian:block-resampling-aa} in
\path{ian/09_gamma_set_blocks_and_forests.tex} holds without
assuming (UAA), with
\[
 \kappa_k=c_\star\qquad\text{for every }k\ge1.
\]
In fact its resampling assertion holds for arbitrary finite
independent blocks without a block-size bound:
\[
 \operatorname{Var}_\mu f\ge v
 \ \Longrightarrow\
 \max_B\mathcal E_B^\mu(f)\ge c_\star(v/d)^{12}.
\]

Corollary~\ref{ian:two-spin-blocks} consequently becomes
\[
 \max_i I_i^\mu(f)\ge c_\star(v/d)^{12}
 \quad\text{for arbitrary coordinate-bias products},
\]
\[
 \max_i I_i^\mu(f)\ge\frac{c_\star}{8}(v/d)^{12}
 \quad\text{for products of blocks of size at most two}.
\]
Both conclusions are hereditary. For Gamma, the Local parameters
are respectively
\[
 (c_L,C_L)=(c_\star,12),\qquad
 (c_L,C_L)=(c_\star/8,12).
\]
Since Local always holds, the four Global parameters may be any
fixed positive constants, for example all equal to one.
\end{proposition}

\begin{proof}
Regard each physical block as one alphabet coordinate.
Every physical monomial of degree at most \(d\) involves at most
\(d\) blocks. Thus \(f\) has interaction degree at most \(d\) on
the product of block alphabets, and remains bounded by one on that
whole product set. Assertion (A-AA) is exactly the required block
resampling bound.

This replaces the dyadic proof of
Lemma~\ref{ian:block-resampling-aa}; its old parameters
\(A_k,M_k\) are unnecessary. The proof of
Corollary~\ref{ian:two-spin-blocks} is unchanged:
\[
 \mathcal E_{\{i\}}^\mu(f)
 =4p_i(1-p_i)I_i^\mu(f)\le I_i^\mu(f),
\]
and for a block of size at most two its existing deterministic
path inequality gives
\[
 \mathcal E_B^\mu(f)\le4\sum_{i\in B}I_i^\mu(f)
 \le8\max_{i\in B}I_i^\mu(f).
\]
Pinning preserves independent residual blocks and cannot increase
interaction degree.
\end{proof}

\begin{remark}[Block and forest failures are unchanged]
\label{ian:influence-consequences-failures}
The block-resampling conclusion does not assert a physical-spin
influence bound for larger blocks. Therefore
Proposition~\ref{ian:hard-block-addressing}, which refutes fixed-set
Gamma from block size three onward, remains valid unchanged.
Its large block energies coexist with vanishing physical influences.

Proposition~\ref{ian:forest-addressing} also remains valid unchanged
for every positive finite edge-strength bound. At zero edge bound,
its positive product-law conclusion is now unconditional.
The forest simulation conclusion continues to use the simulation
form supplied by
Corollary~\ref{cor:resolution-product-alphabet-aa}.
\end{remark}

\begin{proposition}[Updated total-interaction-budget parameters]
\label{ian:influence-consequences-budget}
Proposition~\ref{ian:interaction-budget-aa} in
\path{ian/10_gamma_set_high_temperature.tex} holds without
assuming (UAA). If
\[
 \beta\sum_{|A|\ge2}|J_A|\le B,
\]
then
\[
 \max_iI_i^\mu(f)
 \ge c_\star e^{-26B}(v/d)^{12}.
\]
Thus its Local parameters are
\[
 c_L(B)=c_\star e^{-26B},\qquad C_L=12.
\]
The bound is uniform under pinning, needs no field or bias bound,
and does not require the weighted row bound.
The four Global parameters may again all be set to one.
\end{proposition}

\begin{proof}
The existing proof goes through unchanged with \(c_\star\) in place
of the earlier \(\kappa_1\).
For the field-only product law \(\nu\), its density comparison gives
\[
 \operatorname{Var}_\nu f\ge e^{-2B}v,\qquad
 I_i^\mu(f)\ge e^{-2B}I_i^\nu(f).
\]
Hence
\[
 \max_iI_i^\mu(f)
 \ge e^{-2B}c_\star
       \left(\frac{e^{-2B}v}{d}\right)^{12}
 =c_\star e^{-26B}(v/d)^{12}.
\]
Pinning cannot increase the total absolute mass of residual
interactions of size at least two, so the same constant survives.

More generally a product estimate \(cv^\alpha/d^\beta\) gives
\(ce^{-2(\alpha+1)B}v^\alpha/d^\beta\) by this comparison.
The unnecessary uniform-encoding exponent \(24\) would instead
produce \(e^{-50B}\); the direct product theorem retains \(e^{-26B}\).
\end{proof}

\begin{remark}[Summary of the changes in the consuming subsections]
\label{ian:influence-consequences-status}
Subsections~\ref{ian:exchangeable-generalized-influence},
\ref{ian:generalized-influence-checks},
\ref{ian:gamma-set-blocks-and-forests} and
\ref{ian:gamma-set-high-temperature} no longer assume (UAA).
In the first, the product half of
Lemma~\ref{ian:uniform-product-slice-aa} is (P-AA) and the
constants are (EC); the slice and Theorem~B proofs keep their
exponents and structure. In the second, only the Status line and
the hypothesis of Proposition~\ref{ian:cw-generalized-check} change.
In the third, Lemma~\ref{ian:block-resampling-aa} is proved by the
block-alphabet argument above with \(\kappa_k=c_\star\), and
Corollary~\ref{ian:two-spin-blocks} has constants \(c_\star\) and
\(c_\star/8\); the failure results are unchanged. In the fourth,
Proposition~\ref{ian:interaction-budget-aa} has constant
\(c_\star e^{-26B}\), and the Status line remains ``partial''
because the weighted-row-sum-only assertion is unresolved.
Claims of better exponents elsewhere in the repository still carry
their own (UAA) hypothesis; the \((12,12)\) estimate proved here is
not identified with the stronger statement (UAA).
\end{remark}
